% =====================================================================
% Section 03, Phase 1 original datasets
% =====================================================================

\section{Phase 1 \textbullet\ The Seven Original Datasets}
\label{sec:phase1-datasets}

All seven raw files live in \fn{DatasetsGeojson/} at the repo root.
They arrive in the geographic CRS \stat{EPSG:4326} (lon/lat in degrees)
and are reprojected to \stat{EPSG:32636} (UTM Zone 36N, metres) before
any spatial work.

\subsection{Inventory and what each dataset represents}

\rowcolors{2}{darkpanel}{darkbg}
\begin{table}[H]
\centering\small
\renewcommand{\arraystretch}{1.25}
\begin{tabularx}{\linewidth}{@{}
 >{\color{plotorange}\bfseries}l
 >{\color{bodytext}}l
 >{\color{plotyellow}\ttfamily}r
 >{\color{bodytext}}X @{}}
\toprule
\color{plotblue}File &
\color{plotblue}Geometry &
\color{plotblue}Rows &
\color{plotblue}What it represents \\
\midrule
\fn{Cairo Daily Boarding.json} & Point & 1{,}302 & Bus / microbus stops with daily boarding-and-alighting counts split by time-of-day and by formal/informal mode \\
\fn{Public Transport Routes.json} & LineString & 1{,}784 & Route lines with route code, vehicle type, fare, route length, terminal IDs, and trip metadata \\
\fn{Public Transport Terminals.json} & Point & 280 & Physical terminal locations (\emph{mawaqif}), origin/destination of routes \\
\fn{Public Transport Passenger Flow.json} & LineString & 9{,}258 & Road segments with passengers/hour throughput \\
\fn{Public Transport Vehicles Flow.json} & LineString & 5{,}592 & Road segments with vehicles/hour throughput \\
\fn{Public Transport Commercial Speeds 2.json} & LineString & 26{,}154 & Road segments with average bus speed by time-of-day \\
\fn{Population \& Employment Access.json} & Polygon & 1{,}525 & 1\,km$^2$ H3 hexagonal grid with population, jobs, and accessibility scores \\
\bottomrule
\end{tabularx}
\caption{The seven Phase 1 GeoJSON datasets. Total: \stat{45{,}889}~rows, three geometry types.}
\end{table}
\rowcolors{0}{}{}

\subsection{Dataset profile}

A short profile of each file's most important fields and the questions
it underwrites:

\subsubsection{Boarding (\fn{Cairo Daily Boarding.json})}

The largest behavioural dataset in the project. Each row is a stop
(point) with twelve numeric counts split by time-of-day (morning vs
all-day) and by mode (formal-only, informal-only, all). Key fields:

\begin{itemize}
 \item \fn{Stop\_Name}, human-readable stop name (sometimes Arabic).
 \item \fn{Daily\_All\_BA}, \fn{Daily\_Formal\_BA}, \fn{Daily\_Informal\_BA}
, total daily boardings/alightings, by mode share.
 \item \fn{Morning\_*\_BA}, equivalent counts restricted to the morning
 peak.
\end{itemize}

\textbf{Used by:} Q1 (job accessibility vs. alighting), Q3 (corridor
adherence), Q4 (formal vs. informal dominance), Q6 (route supply vs.
demand), Q7 (morning vs. evening symmetry), Q8 (spatial autocorrelation),
Q9 (underserved hexes), Q11 (terminal-zone classification).

\subsubsection{Routes (\fn{Public Transport Routes.json})}

The supply-side description of the network. Each row is a published
route line with vehicle type, fare, length, and origin/destination
terminal IDs. Critical fields:

\begin{itemize}
 \item \fn{Vehicle\_Type} (\strval{microbus}, \strval{bus}, \strval{minibus}, \strval{tomnaya}, \strval{box}).
 \item \fn{Vehicle\_Capacity}, used to derive the formal/informal
 classification (\fn{Transport\_Class}).
 \item \fn{Fare}, \fn{Route\_Length\_km}.
 \item \fn{Origin\_Terminal\_ID}, \fn{Dest\_Terminal\_ID} (foreign keys to terminals).
\end{itemize}

\textbf{Used by:} Q2 (route symmetry), Q4 (formal vs. informal), Q6 (route
supply vs. passenger demand), Q10 (empty-return), Q11 (route counts per
terminal), Q12 (vehicle-route fit). \textit{This is the most-modified
dataset in cleaning, \stat{45.9\%} of fares are null and require KNN
imputation.}

\subsubsection{Terminals (\fn{Public Transport Terminals.json})}

The 280 \emph{mawaqif} that anchor route origins and destinations. Used
as the spatial join target for boarding-near-terminal and flow-near-terminal
buffers (see merges D, F, G). Both English and Arabic names are
preserved (\fn{Terminal\_Name}, \fn{Terminal\_Name\_AR}).

\subsubsection{Passenger Flow / Vehicle Flow}

These are the segment-level flow datasets. Each row is a road-segment
LineString with hourly throughput. Together they enable the
\stat{Empty Return Index} that defines G2 (\fn{1, passengers/vehicles},
clipped to $[0, 1]$).

\subsubsection{Commercial Speeds}

The largest dataset (26{,}154 rows), segment-level average commercial
speed by time-window. Used by Q12 to compare vehicle types against the
real road conditions they operate on.

\subsubsection{Population \& Employment Access}

The 1{,}525-row demand dataset. Each row is a 1\,km\textsuperscript{2}
H3 hexagon with:

\begin{itemize}
 \item \fn{Population\_2018} (raw count), \fn{Pop\_Male} / \fn{Pop\_Female} splits.
 \item \fn{Jobs\_In\_Hex}, formal employment count.
 \item \fn{Jobs\_Accessible\_60min}, number of jobs reachable within
 a 60-minute commute.
 \item \fn{Jobs\_Pct\_Accessible\_60min}, the same as a percent of all
 jobs in Cairo.
\end{itemize}

This is the dataset that lets us connect transport \emph{infrastructure}
to transport \emph{demand}.

\subsection{Coordinate Reference System and units}

\begin{callout}
\textbf{\color{plotblue}Why we reproject to EPSG:32636.}\quad
The raw files arrive in \stat{EPSG:4326} (lon/lat in degrees). Degrees
are useless for distance work, at Cairo's latitude (30\textdegree~N)
one degree of longitude is $\approx$\stat{96\,km} while one degree of
latitude is $\approx$\stat{111\,km}. Every analysis using a buffer (Q3
120-m corridor, Q11 100-m ghost-terminal test, Phase 2 H3 500-m BRT
corridor), a nearest-neighbour join, or an area calculation must be in
metres. \stat{UTM Zone 36N} is the right projected CRS for Egypt, distances
and areas come out correct without manual scaling. Every cleaned CSV
and every downstream analysis assumes metres.
\end{callout}
